\section{Related Work}
\label{sec:related_work}

\noindent\textbf{Receptor models.}
Chemical mass balance (CMB) and positive matrix factorization (PMF) estimate source
contributions from the measured chemical composition of particulate samples
\citep{watson2002receptor,paatero1994positive,hopke2016review}, and reviews document
their sensitivity to source profiles and collinear tracers
\citep{viana2008source,belis2013critical,belis2019european,mircea2020european,hopke2020global}.
Studies of Indian cities report divergent contributions for the same city
\citep{pant2012critical,karagulian2015contributions}.  \citet{henry1992dealing} defines
an eligible space of estimable combinations for CMB with nearly collinear profiles,
PMF rotational ambiguity is explored with dedicated tools
\citep{paatero2002understanding}, and statistical treatments address temporal
correlation \citep{park2001multivariate} and identification
\citep{jin2025identification}.  Our analysis uses the PM\(_{2.5}\) and wind records of
the New Delhi network, so the sources are separated by transport and timing rather than
by particle composition, and the network's gas measurements are not used.

\noindent\textbf{Source-oriented models and inversion.}
Inverse modelling and data assimilation estimate emissions from concentrations with a
prior \citep{tarantola2005inverse,carrassi2018data}, and atmospheric inversion measures what
the observations constrain through the averaging kernel \citep{rodgers2000inverse} and
chooses how far to aggregate the state by balancing aggregation and smoothing error
\citep{turner2015balancing}.  IASA also aggregates, but without a prior: it reports the
reportable grouping of declared sources with the most groups.
Spatio-temporal models for sensor networks
\citep{li2017diffusion,yu2017spatio,iyer2022modeling,bhardwaj2025comprehensive} target
forecasting and interpolation rather than attribution.

\noindent\textbf{Estimability and collinearity.}
Estimability \citep{searle1971linear}, identification
\citep{rothenberg1971identification,manski2003partial} and collinearity diagnostics
\citep{fox1992generalized,belsley1980regression} do not return a partition.  The
partition of Theorem~\ref{thm:groupings} is the component decomposition of the column
matroid \citep{krogdahl1977dependence,oxley2011matroid}, and the separation is the
sine of a principal angle \citep{bjorck1973numerical}.  Appendix~\ref{app:extended_related_work}
gives a longer discussion.
